\documentclass[a4paper,12pt]{article}

%% ── All standard packages available on Overleaf ──────────────────────────────
\usepackage{lmodern}
\usepackage[T1]{fontenc}
\usepackage[utf8]{inputenc}
\usepackage{geometry}
\geometry{a4paper, top=2.5cm, bottom=2.5cm, left=3cm, right=2.5cm}
\usepackage{graphicx}
\usepackage{amsmath}
\usepackage{makecell}
\usepackage{booktabs}
\usepackage{array}
\usepackage{enumitem}
\usepackage{setspace}
\usepackage{fancyhdr}
\usepackage{titlesec}
\usepackage[backend=biber,style=numeric,sorting=none]{biblatex}
\addbibresource{sample.bib}

%% ── Page style ───────────────────────────────────────────────────────────────
\pagestyle{fancy}
\fancyhf{}
\fancyhead[L]{\small\nouppercase{\leftmark}}
\fancyfoot[C]{\thepage}
\renewcommand{\headrulewidth}{0.4pt}
\setlength{\headheight}{15pt}

%% ── Custom list environments (symbols and abbreviations) ─────────────────────
\newenvironment{symbols}{%
  \begin{list}{}{%
    \setlength{\labelwidth}{3.8cm}%
    \setlength{\leftmargin}{4.4cm}%
    \setlength{\labelsep}{0.5cm}%
    \setlength{\itemsep}{3pt}%
    \renewcommand{\makelabel}[1]{##1\hfil}%
  }%
}{%
  \end{list}%
}

\newenvironment{abbrv}{%
  \begin{list}{}{%
    \setlength{\labelwidth}{2.5cm}%
    \setlength{\leftmargin}{3.2cm}%
    \setlength{\labelsep}{0.5cm}%
    \setlength{\itemsep}{3pt}%
    \renewcommand{\makelabel}[1]{\textbf{##1}\hfil}%
  }%
}{%
  \end{list}%
}

%% ─────────────────────────────────────────────────────────────────────────────
\begin{document}

%% ── TITLE PAGE ───────────────────────────────────────────────────────────────
\begin{titlepage}
  \centering
  \vspace*{1.5cm}
  {\large\bfseries International Burch University\\[0.3cm]}
  {\large Faculty of Engineering\\[0.1cm]}
  {\large Department of Computer Engineering\\[2.5cm]}

  \rule{\linewidth}{1.2pt}\\[0.5cm]
  {\LARGE\bfseries Brain Tumor Segmentation\\[0.3cm]
  Using Deep Learning\\[0.5cm]}
  \rule{\linewidth}{1.2pt}\\[1.5cm]

  {\large Bachelor of Science Thesis\\[2.5cm]}

  \begin{tabular}{@{}ll@{}}
    \textbf{Authors:}    & Salem R S Aker \quad 122200157 \\[0.2cm]
                         & Tarek ALibrahim \quad 122200002 \\[0.2cm]
                         & Faysal Dabbagh \quad 122200015 \\[1cm]
    \textbf{Supervisor:} & Elena Battini Sönmez \\
  \end{tabular}

  \vfill
  {\large \today}
\end{titlepage}

%% ── ABSTRACT ─────────────────────────────────────────────────────────────────
\newpage
\thispagestyle{plain}
\addcontentsline{toc}{section}{Abstract}
\section*{Abstract}
\noindent
Brain tumor segmentation in magnetic resonance imaging (MRI) is an important task for supporting
clinical diagnosis, treatment planning, and tumor monitoring. This study develops and evaluates a
multi-dataset deep learning pipeline for binary brain tumor segmentation. The primary benchmark
dataset is the Figshare Brain Tumor Dataset, which contains 3,064 T1-weighted contrast-enhanced
MRI images and corresponding binary tumor masks from 233 patients. A second, larger dataset,
BraTS2020, comprising approximately 57,000 two-dimensional FLAIR slices extracted from 369
glioma patients, is also incorporated to assess pipeline generalizability.

Seven experimental configurations are systematically evaluated on the Figshare dataset to isolate
the contributions of loss function design, preprocessing strategy, encoder architecture, and
skip-connection structure. All configurations are based on U-Net-style encoder--decoder
architectures and are evaluated under five-fold cross-validation using an image-level splitting
strategy that matches the publicly available reference benchmark. Performance is measured using
the Dice Similarity Coefficient (DSC) and Intersection-over-Union (IoU), with globally pooled
(micro-averaged) reductions to handle class imbalance.

A preliminary result for fold~1 of the U-Net++ EfficientNet-B4 DiceBCE configuration yields a
test Dice of 0.8706, which is consistent with and marginally above the five-fold reference
benchmark of 0.8608 for this architecture. Full cross-validation results for all seven
configurations are in progress.

%% ── TABLE OF CONTENTS ────────────────────────────────────────────────────────
\newpage
\addcontentsline{toc}{section}{Contents}
\tableofcontents

\newpage
\addcontentsline{toc}{section}{List of Figures}
\listoffigures

\newpage
\addcontentsline{toc}{section}{List of Tables}
\listoftables

%% ── LIST OF SYMBOLS ──────────────────────────────────────────────────────────
\newpage
\addcontentsline{toc}{section}{List of Symbols}
\section*{List of Symbols}
\begin{symbols}
    \item[$I(x,y)$]               Intensity value of an MRI image at pixel location $(x,y)$
    \item[$P$]                    Predicted binary tumor segmentation mask
    \item[$G$]                    Ground-truth binary tumor segmentation mask
    \item[$P \cap G$]             Overlapping region between the predicted mask and the ground-truth mask
    \item[$P \cup G$]             Union region between the predicted mask and the ground-truth mask
    \item[$S_e$]                  Feature map extracted from the encoder path of the segmentation model
    \item[$S_d$]                  Feature map reconstructed in the decoder path of the segmentation model
    \item[$S_C$]                  Concatenated feature map formed using skip connections
    \item[$p(x,y)$]               Predicted tumor probability at pixel location $(x,y)$
    \item[$\mathcal{L}_{Dice}$]   Dice loss function
    \item[$\mathcal{L}_{BCE}$]    Binary cross-entropy loss function
    \item[$\mathcal{L}_{DiceBCE}$] Combined Dice and binary cross-entropy loss function
    \item[$DSC$]                  Dice Similarity Coefficient
    \item[$IoU$]                  Intersection-over-Union
\end{symbols}

%% ── LIST OF ABBREVIATIONS ────────────────────────────────────────────────────
\newpage
\addcontentsline{toc}{section}{List of Abbreviations}
\section*{List of Abbreviations}
\begin{abbrv}
    \item[BCE]      Binary Cross-Entropy
    \item[BraTS]    Brain Tumor Segmentation (benchmark challenge)
    \item[CLAHE]    Contrast Limited Adaptive Histogram Equalization
    \item[CNN]      Convolutional Neural Network
    \item[CT]       Computed Tomography
    \item[CV]       Cross-Validation
    \item[DSC]      Dice Similarity Coefficient
    \item[FLAIR]    Fluid Attenuated Inversion Recovery
    \item[GPU]      Graphics Processing Unit
    \item[IoU]      Intersection-over-Union
    \item[MRI]      Magnetic Resonance Imaging
    \item[PET]      Positron Emission Tomography
    \item[SMP]      Segmentation Models PyTorch
    \item[U-Net]    U-shaped Encoder--Decoder Convolutional Neural Network
    \item[U-Net++]  Nested U-Net Architecture
\end{abbrv}

\newpage
\pagenumbering{arabic}
\setcounter{page}{1}

%% ─────────────────────────────────────────────────────────────────────────────
%% SECTIONS
%% ─────────────────────────────────────────────────────────────────────────────

\section{Introduction}
\label{sec:introduction}

The brain is the central organ responsible for controlling and regulating essential bodily
functions, including cognition, motor control, and sensory processing. The abnormal and
uncontrolled proliferation of cells within brain tissue leads to the formation of brain tumors,
which can significantly impair neurological function and pose serious health risks
\cite{louis2021who}. Brain tumors are broadly classified into benign and malignant types based
on their growth behavior, invasiveness, and histopathological characteristics. Benign tumors
typically exhibit slow growth and well-defined boundaries, whereas malignant tumors demonstrate
rapid proliferation, infiltrative growth, and aggressive invasion of surrounding healthy tissue
\cite{who2021tumors}.

Among malignant brain tumors, gliomas represent one of the most common and deadliest categories,
originating from glial cells within the central nervous system. High-grade gliomas, particularly
glioblastoma, are characterized by rapid progression, extensive infiltration, and resistance to
treatment, leading to a poor prognosis. Despite advances in surgery, radiotherapy, and
chemotherapy, the median survival time for patients diagnosed with high-grade gliomas remains
approximately 12--24 months \cite{ostrom2020cbtrus,louis2021who}.

Medical imaging plays a critical role in the diagnosis, treatment planning, and follow-up of
brain tumors. Among available imaging modalities, Magnetic Resonance Imaging (MRI) is the
primary tool used in clinical practice due to its superior soft-tissue contrast and absence of
ionizing radiation \cite{shukla2017advanced}. MRI enables the visualization of tumor
sub-regions, such as the enhancing core, necrotic tissue, and surrounding edema, which are
essential for accurate diagnosis and therapeutic decision-making.

However, precise delineation of tumor boundaries from MRI scans is typically performed manually
by medical experts, making the process time-consuming, labor-intensive, and subject to
inter-observer variability \cite{amin2025dual}. Automated brain tumor segmentation remains
computationally challenging due to several intrinsic factors, including irregular tumor shapes,
fuzzy growth patterns, image noise artifacts, partial volume effects, and severe class imbalance
between tumor pixels and healthy tissue \cite{tohka2014partial}.

To address these challenges, recent research has focused on developing deep learning systems for
automated tumor segmentation. Among these methods, encoder--decoder architectures have become
widely adopted in medical image segmentation because they can learn high-level visual features
while preserving spatial details through skip connections. U-Net-based models are especially
suitable for this task because they produce pixel-wise predictions and can be trained directly
using paired MRI images and ground-truth tumor masks.

This project develops and evaluates a multi-dataset binary brain tumor segmentation pipeline.
Seven experimental configurations are compared across two datasets---the single-modality
Figshare Brain Tumor Dataset and the multi-patient multi-modal BraTS2020 dataset---to assess how
loss function design, preprocessing strategy, encoder architecture, and skip-connection structure
affect segmentation performance:

\begin{itemize}
    \item Five configurations sharing the same U-Net architecture and ResNet-34 encoder but
          differing in their loss function: Dice, BCE, DiceBCE, DiceFocal, and Lovász.
    \item One configuration adding CLAHE contrast enhancement during preprocessing to the
          DiceBCE baseline (CLAHE DiceBCE).
    \item One configuration replacing the baseline U-Net encoder and skip-connection design with
          a U-Net++ architecture and an EfficientNet-B4 encoder (U-Net++ EfficientNet-B4 DiceBCE).
\end{itemize}

The remainder of this thesis is organized as follows. Section~\ref{sec:related_works} reviews
existing studies on brain tumor segmentation. Section~\ref{sec:methodology} describes the
datasets, preprocessing steps, and segmentation models. Section~\ref{sec:Experiments} presents
the experimental setup and five-fold cross-validation strategy. Section~\ref{EvaluationMetrics}
defines the evaluation metrics. Section~\ref{sec:results} presents and discusses the
experimental results, and Section~\ref{sec:conclusion} concludes the study.


\section{Related Work}
\label{sec:related_works}

Brain tumor segmentation has been extensively studied, with earlier methods relying on
traditional machine learning and clustering techniques while more recent approaches leverage
deep learning architectures. A comparison of representative methods, organized by algorithm,
dataset, and reported Dice performance, is presented in Table~\ref{tab:related_works_comparison}.

Zhang et al.\ \cite{zhang2019brain} employed a hybrid clustering approach based on K-means,
achieving a Dice score of approximately 85\% on a private MRI dataset. As deep learning gained
prominence, Saouli et al.\ \cite{saouli2018fully} employed an Incremental Deep Neural Network
on the BRATS dataset, improving performance to 88.6\%.

Significant advances were made by Kamnitsas et al.\ \cite{kamnitsas2017efficient} with the 3D
CNN (DeepMedic) architecture, which effectively captures volumetric context and yields a Dice
score of 90.1\%. Building on the U-Net architecture, Allah et al.\ \cite{allah2023edge} proposed
the Edge U-Net, achieving 91.3\%, while Ruba et al.\ \cite{ruba2023brain} introduced the
JGate-AttResU-Net, pushing accuracy to 92.1\%.

These results demonstrate that improvements in network architecture, preprocessing, and loss
function design can produce consistent gains in segmentation accuracy. Based on this direction,
the current study develops a systematic pipeline that ablates five loss functions on a common
U-Net backbone, evaluates CLAHE preprocessing as an independent variable, and tests whether a
stronger architecture (U-Net++ with EfficientNet-B4) provides consistent gains over the baseline.

\begin{table}[htbp]
    \centering
    \caption{Comparison of existing brain tumor segmentation methods and the current project approach.}
    \label{tab:related_works_comparison}
    \renewcommand{\arraystretch}{1.3}
    \resizebox{\textwidth}{!}{
        \begin{tabular}{|l|l|l|l|}
        \hline
        \textbf{Paper / Study} & \textbf{Algorithm} & \textbf{Dataset} & \textbf{Performance} \\ \hline

        Zhang et al.\ (2019) \cite{zhang2019brain}
        & \makecell{Hybrid Clustering \\ (K-means)}
        & \makecell{Private MRI \\ Dataset}
        & Dice: $\sim$85\% \\ \hline

        Saouli et al.\ (2018) \cite{saouli2018fully}
        & \makecell{Incremental \\ DNN}
        & \makecell{BRATS \\ Dataset}
        & Dice: 88.6\% \\ \hline

        Kamnitsas et al.\ (2017) \cite{kamnitsas2017efficient}
        & \makecell{3D CNN \\ (DeepMedic)}
        & \makecell{BRATS \\ Dataset}
        & Dice: 90.1\% \\ \hline

        Allah et al.\ (2023) \cite{allah2023edge}
        & Edge U-Net
        & \makecell{BRATS \\ Dataset}
        & Dice: 91.3\% \\ \hline

        Ruba et al.\ (2023) \cite{ruba2023brain}
        & JGate-AttResU-Net
        & \makecell{BRATS \\ Dataset}
        & Dice: 92.1\% \\ \hline

        Current Study
        & \makecell{U-Net (ResNet-34, 5 losses),\\
                    CLAHE U-Net (DiceBCE),\\
                    U-Net++ (EfficientNet-B4)}
        & \makecell{Figshare + \\ BraTS2020}
        & \makecell{Five-fold \\ cross-validation \\ (in progress)} \\ \hline

        \end{tabular}
    }
\end{table}

\begin{figure}[htbp]
    \centering
    \includegraphics[width=0.85\textwidth]{u-net-architecture.png}
    \caption{Overview of the U-Net architecture commonly used for biomedical image segmentation \cite{ronneberger2015unet}.}
    \label{fig:unet_architecture}
\end{figure}


\newpage

\section{Methodology}
\label{sec:methodology}

This section describes the datasets, preprocessing steps, segmentation models, training strategy,
and evaluation protocol used in this study. The project develops an independent deep learning
pipeline for binary brain tumor segmentation that is evaluated across two complementary
datasets and seven experimental configurations.

The proposed workflow begins by extracting MRI images and tumor masks from the original raw
files, resizing and normalizing all image--mask pairs, optionally applying contrast enhancement,
applying data augmentation during training, and training segmentation models under five-fold
cross-validation. The seven evaluated configurations systematically vary the loss function (Dice,
BCE, DiceBCE, DiceFocal, Lovász), preprocessing (with and without CLAHE), and architecture
(U-Net with ResNet-34 vs.\ U-Net++ with EfficientNet-B4).


\subsection{Datasets}
\label{sub:datasets}

\subsubsection{Figshare Brain Tumor Dataset}
\label{sub:figshare}

The primary dataset used in this study is the Figshare Brain Tumor Dataset
\cite{cheng2017figshare}, a publicly available MRI collection prepared for brain tumor analysis.
The dataset contains 3,064 T1-weighted contrast-enhanced MRI images collected from 233 patients.
Each sample is provided in \texttt{.mat} file format and includes the raw MRI image, the
corresponding binary tumor mask, a tumor class label (meningioma, glioma, or pituitary), and a
patient identifier.

The task is formulated as binary semantic segmentation: each pixel is classified as either tumor
(foreground) or background. The \texttt{.mat} files are converted into paired PNG image--mask
samples. Each MRI image is converted from grayscale to three-channel format to match the input
requirements of ImageNet-pretrained encoder weights. All images and masks are resized to
$256 \times 256$ pixels. Pixel intensity values are normalized using ImageNet mean and standard
deviation statistics, and tumor masks are binarized (foreground = 1, background = 0) before
training.

\subsubsection{BraTS2020 Dataset}
\label{sub:brats2020}

The BraTS2020 dataset \cite{bakas2018identifying} is used as a second evaluation environment to
assess whether the segmentation pipeline generalizes beyond the single-modality Figshare setting.
The dataset consists of pre-operative multi-parametric MRI scans from 369 glioma patients,
provided as three-dimensional volumes. For compatibility with the two-dimensional pipeline
developed in this study, the FLAIR modality channel is extracted and each three-dimensional
volume is decomposed into individual axial slices. Slices containing no tumor tissue are filtered
out, yielding approximately 57,000 tumor-bearing two-dimensional slices.

Three tumor sub-region channels (necrotic core, peritumoral edema, and enhancing tumor) are
merged into a single binary whole-tumor mask at preprocessing time. Slices and masks are
resized to $256 \times 256$ pixels using the same normalization and binarization procedure
applied to the Figshare dataset. Because all BraTS2020 patients are diagnosed with glioma,
the tumor class label is uniform, and patient-disjoint (patient-level) cross-validation splits
are used to prevent data leakage.

\begin{table}[htbp]
    \centering
    \caption{Summary of the datasets used in this study.}
    \label{tab:dataset_summary}
    \renewcommand{\arraystretch}{1.5}
    \resizebox{\textwidth}{!}{
    \begin{tabular}{|l|l|l|l|l|}
        \hline
        \textbf{Dataset} & \textbf{Samples} & \textbf{Modality} & \textbf{Classes} & \textbf{CV Scheme} \\ \hline

        \makecell[l]{Figshare \cite{cheng2017figshare}}
        & \makecell[l]{3,064 images \\ 233 patients}
        & T1-weighted
        & \makecell[l]{Meningioma \\ Glioma \\ Pituitary}
        & Image-level KFold \\ \hline

        BraTS2020 \cite{bakas2018identifying}
        & \makecell[l]{$\sim$57,000 slices \\ 369 patients}
        & FLAIR
        & Glioma
        & Patient-level GroupKFold \\ \hline

    \end{tabular}
    }
\end{table}

\begin{figure}[h!]
    \centering
    \includegraphics[width=0.45\textwidth]{Figshare.png}
    \caption{Example MRI slice from the Figshare Brain Tumor Dataset used for binary tumor segmentation.}
    \label{fig:dataset_examples}
\end{figure}


\subsection{Architectures}
\label{sub:architectures}

\subsubsection{U-Net with ResNet-34 Encoder}
\label{sub:unet_architecture}

U-Net is an encoder--decoder convolutional neural network architecture widely used for biomedical
image segmentation \cite{ronneberger2015unet}. The encoder path extracts hierarchical image
features by applying convolutional operations and progressive downsampling, while the decoder path
reconstructs the spatial resolution to produce a pixel-wise segmentation mask. Skip connections
between corresponding encoder and decoder levels preserve spatial details that may be lost during
downsampling.

Five experimental configurations (Experiments~01--05 and~06) use U-Net with a ResNet-34 encoder
pre-trained on ImageNet as the feature extraction backbone. The decoder produces a single-channel
output at the same spatial resolution as the input. A sigmoid activation is applied at inference
time to produce a tumor probability map, and the binary mask is obtained by thresholding at~0.5.

\subsubsection{U-Net++ with EfficientNet-B4 Encoder}
\label{sub:unetpp_effb4}

U-Net++ is an extension of the original U-Net architecture that redesigns the skip connections
using nested and densely connected convolutional pathways \cite{zhou2019unetpp}. This structure
reduces the semantic gap between encoder and decoder feature maps by inserting intermediate
nodes along the skip paths, enabling the model to combine multi-scale features more effectively
than standard skip connections.

Experiment~07 combines U-Net++ with an EfficientNet-B4 encoder pre-trained on ImageNet.
EfficientNet-B4 employs compound scaling to balance network depth, width, and resolution,
providing richer feature representations than a ResNet-34 encoder. This configuration evaluates
whether the stronger encoder and nested skip-connection structure improve tumor segmentation
performance relative to the baseline. The batch size is reduced from 8 to 6 to accommodate the
larger model within GPU memory constraints.


\subsection{Loss Functions}
\label{sub:losses}

All seven configurations are trained with the same optimizer and scheduler but differ in the
loss function applied during training. The following loss functions are evaluated:

\subsubsection{Binary Cross-Entropy (BCE)}
Binary Cross-Entropy applies standard pixel-wise classification between tumor and background.
It treats each pixel independently and is susceptible to class imbalance when the tumor region
is small relative to the overall image area:
\begin{equation}
    \mathcal{L}_{BCE} = -\frac{1}{N}\sum_{i=1}^{N}
        \bigl[y_i \log \hat{y}_i + (1 - y_i) \log(1 - \hat{y}_i)\bigr]
\end{equation}
where $y_i$ is the ground-truth label and $\hat{y}_i$ is the predicted probability for pixel~$i$.

\subsubsection{Dice Loss}
Dice loss directly optimizes the Dice Similarity Coefficient, the primary evaluation metric:
\begin{equation}
    \mathcal{L}_{Dice} = 1 - \frac{2\,|P \cap G| + \epsilon}{|P| + |G| + \epsilon}
\end{equation}
where $\epsilon$ is a smoothing constant that prevents division by zero. Dice loss is less
sensitive to class imbalance than BCE because it operates on set-level overlap rather than
individual pixel predictions.

\subsubsection{Combined Dice-BCE (DiceBCE)}
The combined DiceBCE loss function adds BCE and Dice loss with equal weights:
\begin{equation}
    \mathcal{L}_{DiceBCE} = \mathcal{L}_{Dice} + \mathcal{L}_{BCE}
\end{equation}
By combining both terms, the model is encouraged to produce accurate pixel-level classification
(BCE) while also maximizing the spatial overlap between the predicted and ground-truth tumor
regions (Dice). This is the primary baseline configuration used in this study.

\subsubsection{Combined Dice-Focal (DiceFocal)}
Focal loss \cite{lin2017focal} is a variant of cross-entropy that down-weights easy, well-classified
pixels and focuses training on hard examples:
\begin{equation}
    \mathcal{L}_{Focal} = -\alpha\,{(1 - \hat{y})}^{\gamma}\,\log(\hat{y})
\end{equation}
with $\alpha = 0.25$ and $\gamma = 2.0$. The DiceFocal configuration pairs Focal loss with Dice
loss to combine hard-sample weighting with global overlap optimization.

\subsubsection{Lovász Loss}
The Lovász hinge loss \cite{berman2018lovasz} is a surrogate for the IoU metric derived via the
Lovász extension of submodular functions. It directly optimizes the IoU score at the set level,
offering a differentiable loss that is theoretically better aligned with the evaluation metric
than either BCE or Dice loss.

\subsubsection{CLAHE Preprocessing}
Experiment~06 applies Contrast Limited Adaptive Histogram Equalization (CLAHE) to MRI images
before training. MRI scans may contain low-contrast regions where tumor boundaries are difficult
to distinguish from surrounding tissue. CLAHE improves local intensity differences by performing
histogram equalization on small, overlapping image tiles while limiting noise amplification
through a user-defined contrast limit. The model architecture, loss function (DiceBCE), and all
other training parameters remain identical to Experiment~03, isolating the effect of CLAHE.


\section{Experimental Setup}
\label{sec:Experiments}

This section describes the experimental setup used to train and evaluate the seven segmentation
configurations. All experiments are conducted under the same data preparation procedure, image
resolution, and evaluation protocol to ensure fair and controlled comparisons.

The original MRI images and tumor masks are extracted from the dataset's raw files and converted
into paired PNG image--mask samples. Each image and its corresponding mask are resized to
$256 \times 256$ pixels. Grayscale MRI images are converted to three-channel format to match
the input requirements of the ImageNet-pretrained encoders. Intensity values are normalized using
ImageNet statistics, and the tumor masks are binarized before training.

Five-fold cross-validation is applied using an image-level KFold splitting strategy on the
Figshare dataset, matching the publicly available reference benchmark and enabling direct
comparison of results. During each fold, four partitions are used for training (with a 15\%
validation subset drawn from the training pool) and one partition is reserved for testing.
This process is repeated five times so that each fold serves once as the test set. The final
performance is reported as the mean and standard deviation of the metric across the five test
folds. For BraTS2020, a patient-disjoint GroupKFold strategy is used to guarantee that no
patient's slices appear in both training and test partitions within the same fold.


\subsection{Experimental Configurations}
\label{sub:experimental_configurations}

Table~\ref{tab:configurations} summarizes the seven configurations evaluated in this study.
Experiments~01--05 form a controlled loss-function ablation with the architecture held constant.
Experiment~06 introduces CLAHE preprocessing while keeping the architecture and loss identical
to Experiment~03. Experiment~07 upgrades the encoder and skip-connection design.

\begin{table}[htbp]
\centering
\caption{Summary of the seven experimental configurations. All share the same U-Net/ResNet-34
architecture except Experiment~07. All use the DiceBCE loss except where noted.}
\label{tab:configurations}
\renewcommand{\arraystretch}{1.3}
\resizebox{\textwidth}{!}{
\begin{tabular}{|c|l|l|l|l|}
\hline
\textbf{Exp.} & \textbf{Name} & \textbf{Architecture} & \textbf{Loss} & \textbf{Preprocessing} \\ \hline
01 & Dice             & U-Net / ResNet-34           & Dice       & None (baseline) \\ \hline
02 & BCE              & U-Net / ResNet-34           & BCE        & None \\ \hline
03 & DiceBCE          & U-Net / ResNet-34           & Dice + BCE & None \\ \hline
04 & DiceFocal        & U-Net / ResNet-34           & Dice + Focal & None \\ \hline
05 & Lovász           & U-Net / ResNet-34           & Lovász     & None \\ \hline
06 & CLAHE DiceBCE    & U-Net / ResNet-34           & Dice + BCE & CLAHE \\ \hline
07 & U-Net++ EffB4    & U-Net++ / EfficientNet-B4   & Dice + BCE & None \\ \hline
\end{tabular}
}
\end{table}


\subsection{Training Parameters}
\label{sub:training_parameters}

All models are trained from image--mask pairs using the hyperparameters summarized in
Table~\ref{tab:training_params}. The Adam optimizer is used with an initial learning rate of
$10^{-4}$. The learning rate is reduced during training using a ReduceLROnPlateau scheduler,
which decreases the learning rate by a factor of 0.1 when the validation loss does not improve
for five consecutive epochs, with a minimum learning rate of $10^{-7}$. Training is stopped
early if the validation Dice score does not improve for 15 consecutive epochs. The best
checkpoint across all epochs is selected based on the highest validation Dice score. A fixed
random seed of 42 is used across all experiments to ensure reproducibility.

Data augmentation is applied during training to improve generalization and reduce overfitting,
following the augmentation strategy used in the publicly available Figshare reference benchmark.
The augmentation pipeline includes horizontal flipping, random affine transformation (scale,
rotate, translate), elastic distortion, brightness and contrast variation, Gaussian blur, and
Gaussian noise.

\begin{table}[htbp]
    \centering
    \caption{Training hyperparameters shared across all experimental configurations.}
    \label{tab:training_params}
    \renewcommand{\arraystretch}{1.3}
    \begin{tabular}{|l|l|}
        \hline
        \textbf{Hyperparameter} & \textbf{Value} \\ \hline
        Image size               & $256 \times 256$ \\ \hline
        Batch size               & 8 \;(6 for Experiment~07) \\ \hline
        Optimizer                & Adam \\ \hline
        Initial learning rate    & $10^{-4}$ \\ \hline
        LR scheduler             & ReduceLROnPlateau \\ \hline
        Scheduler factor         & 0.1 \\ \hline
        Scheduler patience       & 5 epochs (monitor: validation loss) \\ \hline
        Minimum learning rate    & $10^{-7}$ \\ \hline
        Maximum epochs           & 100 \\ \hline
        Early stopping patience  & 15 epochs (monitor: validation Dice) \\ \hline
        Segmentation threshold   & 0.5 \\ \hline
        Random seed              & 42 \\ \hline
        Encoder weights          & ImageNet pre-trained \\ \hline
    \end{tabular}
\end{table}


\section{Evaluation Metrics}
\label{EvaluationMetrics}

The performance of the segmentation models is evaluated by comparing the predicted tumor masks
with the corresponding ground-truth masks. Since the task is binary tumor segmentation, each
pixel is classified as either tumor or background. The two primary evaluation metrics used in
this study are the Dice Similarity Coefficient (DSC) and the Intersection-over-Union (IoU).

The Dice Similarity Coefficient is defined as:

\begin{equation}
    DSC = \frac{2 |P \cap G|}{|P| + |G|}
\end{equation}

where $P$ is the predicted tumor mask and $G$ is the ground-truth tumor mask.
The term $|P \cap G|$ represents the overlapping region between the predicted and actual tumor
areas. A Dice score close to 1 indicates strong spatial overlap and accurate segmentation, while
a score close to 0 indicates poor agreement with the ground truth.

The Intersection-over-Union is defined as:

\begin{equation}
    IoU = \frac{|P \cap G|}{|P \cup G|}
\end{equation}

where $|P \cup G|$ is the union of the predicted and ground-truth masks. IoU provides a
stricter measure of segmentation quality than DSC because the union penalizes false positives
and false negatives more strongly.

Both metrics are computed using globally pooled pixel counts across all images in a test fold
(micro-averaged reduction). This pooling strategy avoids the instability of averaging per-image
scores when the ground-truth mask is very small, which can occur for slices that contain only a
few tumor pixels.

In addition to the primary metrics, the pipeline also records per-image sensitivity
(recall = TP / (TP + FN)) and precision (TP / (TP + FP)), as well as predicted-mask and
ground-truth mask area ratios, to support per-class and per-patient analysis. These are written
to a per-image CSV file alongside the primary Dice and IoU values.


\section{Results \& Discussion}
\label{sec:results}

This section presents the experimental results obtained from the five-fold cross-validation
evaluation on the Figshare Brain Tumor Dataset for the seven experimental configurations.
Performance is measured using the Dice Similarity Coefficient and Intersection-over-Union.

\subsection{Reference Benchmark Comparison}
\label{sub:reference_benchmark}

Table~\ref{tab:reference_results} lists the five-fold mean Dice scores reported in the publicly
available reference benchmark for the Figshare dataset. These values serve as reproduction
targets for each corresponding configuration in this study.

\begin{table}[htbp]
\centering
\caption{Published five-fold reference benchmark results on the Figshare dataset (image-level
KFold, batch~8, Adam $lr=10^{-4}$, ReduceLROnPlateau, max\_epochs=100, seed=42). Experiment~07
uses batch~6.}
\label{tab:reference_results}
\renewcommand{\arraystretch}{1.3}
\begin{tabular}{|c|l|c|}
\hline
\textbf{Exp.} & \textbf{Configuration} & \textbf{Reference Mean Dice} \\ \hline
01 & Dice             & 0.8426 \\ \hline
02 & BCE              & 0.8485 \\ \hline
03 & DiceBCE          & 0.8541 \\ \hline
04 & DiceFocal        & 0.8509 \\ \hline
05 & Lovász           & 0.8418 \\ \hline
06 & CLAHE DiceBCE    & 0.8501 \\ \hline
07 & U-Net++ EffB4 DiceBCE & \textbf{0.8608} \\ \hline
\end{tabular}
\end{table}

\subsection{Fold-wise Experimental Results}
\label{sub:fold_results}

Table~\ref{tab:fold_results} presents the fold-wise Dice and IoU results obtained by the
pipeline for each experimental configuration. Cells marked with `---' indicate that the
corresponding training run has not yet completed.

\begin{table}[htbp]
\centering
\caption{Fold-wise Dice / IoU results for the seven segmentation configurations on the Figshare
dataset (image-level KFold). Cells marked `---' are pending.}
\label{tab:fold_results}
\renewcommand{\arraystretch}{1.3}
\resizebox{\textwidth}{!}{
\begin{tabular}{|c|cc|cc|cc|cc|cc|cc|cc|}
\hline
\multirow{2}{*}{\textbf{Fold}}
  & \multicolumn{2}{c|}{\textbf{Exp.01 Dice}}
  & \multicolumn{2}{c|}{\textbf{Exp.02 BCE}}
  & \multicolumn{2}{c|}{\textbf{Exp.03 DiceBCE}}
  & \multicolumn{2}{c|}{\textbf{Exp.04 DiceFocal}}
  & \multicolumn{2}{c|}{\textbf{Exp.05 Lovász}}
  & \multicolumn{2}{c|}{\textbf{Exp.06 CLAHE}}
  & \multicolumn{2}{c|}{\textbf{Exp.07 U-Net++}} \\
& \textbf{Dice} & \textbf{IoU}
& \textbf{Dice} & \textbf{IoU}
& \textbf{Dice} & \textbf{IoU}
& \textbf{Dice} & \textbf{IoU}
& \textbf{Dice} & \textbf{IoU}
& \textbf{Dice} & \textbf{IoU}
& \textbf{Dice} & \textbf{IoU} \\ \hline

1 & --- & --- & --- & --- & --- & --- & --- & --- & --- & --- & --- & --- & 0.8706 & 0.7708 \\ \hline
2 & --- & --- & --- & --- & --- & --- & --- & --- & --- & --- & --- & --- & ---    & ---    \\ \hline
3 & --- & --- & --- & --- & --- & --- & --- & --- & --- & --- & --- & --- & ---    & ---    \\ \hline
4 & --- & --- & --- & --- & --- & --- & --- & --- & --- & --- & --- & --- & ---    & ---    \\ \hline
5 & --- & --- & --- & --- & --- & --- & --- & --- & --- & --- & --- & --- & ---    & ---    \\ \hline

\end{tabular}
}
\end{table}


\subsection{Average Performance Comparison}
\label{sub:average_results}

Table~\ref{tab:average_results} will be populated with five-fold mean Dice and IoU values once
all training runs have completed. The reference benchmarks from Table~\ref{tab:reference_results}
are reproduced here for direct comparison.

\begin{table}[htbp]
\centering
\caption{Mean Dice and IoU across five folds: our results vs.\ the reference benchmark.
  Our results are pending except where noted.}
\label{tab:average_results}
\renewcommand{\arraystretch}{1.3}
\begin{tabular}{|c|l|cc|cc|}
\hline
\textbf{Exp.} & \textbf{Configuration}
  & \textbf{Ours: Mean Dice} & \textbf{Ours: Mean IoU}
  & \textbf{Ref. Dice} & \textbf{Ref. IoU} \\ \hline
01 & Dice               & --- & --- & 0.8426 & --- \\ \hline
02 & BCE                & --- & --- & 0.8485 & --- \\ \hline
03 & DiceBCE            & --- & --- & 0.8541 & --- \\ \hline
04 & DiceFocal          & --- & --- & 0.8509 & --- \\ \hline
05 & Lovász             & --- & --- & 0.8418 & --- \\ \hline
06 & CLAHE DiceBCE      & --- & --- & 0.8501 & --- \\ \hline
07 & U-Net++ EffB4      & --- & --- & \textbf{0.8608} & --- \\ \hline
\end{tabular}
\end{table}


\subsection{Discussion}
\label{sub:discussion}

The only completed result currently available is fold~1 of Experiment~07 (U-Net++
EfficientNet-B4 DiceBCE), which achieved a test Dice of 0.8706 and a test IoU of 0.7708.
This result exceeds the five-fold reference mean of 0.8608 for this architecture, which is
consistent with expected fold-level variance: individual folds differ from the five-fold mean
depending on the relative difficulty of each test partition.

The training curve for fold~1 of Experiment~07 shows healthy convergence. The validation Dice
improved steadily from 0.732 at epoch~0 to a peak of approximately 0.863 before early stopping
was triggered at epoch~61. The learning rate decayed through the ReduceLROnPlateau schedule,
dropping from $10^{-4}$ to $10^{-5}$ at epoch~30, to $10^{-6}$ at epoch~51, and to $10^{-7}$
at epoch~57. This decay pattern is consistent with the model converging and the optimizer making
progressively finer adjustments, suggesting the pipeline is correctly implemented.

Once all seven configurations have been evaluated, the following comparisons will be of primary
interest:

\begin{itemize}
    \item \textbf{Loss ablation (Exps.\ 01--05):} Isolates the effect of loss function choice
          on segmentation accuracy when the architecture, preprocessing, optimizer, and training
          protocol are held constant. This comparison directly answers which loss function is
          most effective for this dataset.

    \item \textbf{CLAHE effect (Exp.\ 03 vs.\ 06):} Isolates the effect of local contrast
          enhancement during preprocessing. All other factors are identical, so any difference
          in mean Dice is attributable solely to CLAHE.

    \item \textbf{Architecture upgrade (Exp.\ 03 vs.\ 07):} Isolates the combined effect of
          using a stronger encoder (EfficientNet-B4 vs.\ ResNet-34) and a more expressive
          skip-connection design (U-Net++ vs.\ U-Net). This comparison quantifies the
          performance gain from a more capable architecture relative to the baseline.

    \item \textbf{Best preprocessing vs.\ best architecture:} Once mean Dice values are
          available, the CLAHE gain (Exp.\ 06 $-$ Exp.\ 03) can be compared with the
          architecture gain (Exp.\ 07 $-$ Exp.\ 03) to assess which intervention yields a
          larger improvement.
\end{itemize}

The reference benchmark (Table~\ref{tab:reference_results}) indicates that, across all seven
configurations, the performance range is narrow: from 0.8418 (Lovász, Exp.\ 05) to 0.8608
(U-Net++ EfficientNet-B4, Exp.\ 07). The preliminary fold~1 result of 0.8706 for Experiment~07
is consistent with this range and supports the hypothesis that the stronger encoder and nested
skip connections offer a meaningful but not dramatic improvement over the ResNet-34 U-Net
baseline.


\section{Conclusion}
\label{sec:conclusion}

This study presents a multi-dataset deep learning pipeline for binary brain tumor segmentation.
The pipeline converts raw MRI data from two datasets---the Figshare Brain Tumor Dataset
(3,064 T1-weighted images, 233 patients) and the BraTS2020 dataset (${\sim}57{,}000$ FLAIR
slices, 369 patients)---into standardized $256 \times 256$ PNG image--mask pairs, applies
normalization and data augmentation, and evaluates models through five-fold cross-validation.
Seven experimental configurations are compared on the Figshare dataset to systematically
isolate the contributions of loss function choice (Dice, BCE, DiceBCE, DiceFocal, Lovász),
preprocessing (CLAHE), and architecture (U-Net++ with EfficientNet-B4).

All configurations share the same optimizer (Adam, $lr = 10^{-4}$), learning rate schedule
(ReduceLROnPlateau, factor 0.1, patience 5), early stopping criterion (patience 15 epochs on
validation Dice), and random seed (42), providing a controlled basis for comparison. Evaluation
is currently in progress. A preliminary fold~1 result for the U-Net++ EfficientNet-B4 DiceBCE
configuration yields a test Dice of 0.8706 and a test IoU of 0.7708, consistent with the
reference benchmark of 0.8608 for this architecture. The training dynamics for this fold indicate
healthy convergence with appropriate learning rate decay, confirming that the pipeline is
correctly implemented.

The remaining four folds for Experiment~07 and all five folds for Experiments~01--06 are
currently running. Once all evaluations are complete, five-fold mean Dice and IoU will be
reported for each configuration and compared against the published reference benchmarks.
The comparisons will quantify (i) the relative impact of five different loss functions on a
fixed U-Net architecture, (ii) the isolated benefit of CLAHE preprocessing, and (iii) the
performance gain from upgrading to a U-Net++ backbone with an EfficientNet-B4 encoder.

Evaluation on the BraTS2020 dataset using patient-disjoint cross-validation splits will follow,
providing a more rigorous measure of generalization than image-level splits and enabling
assessment of how well the segmentation pipeline transfers to a larger, more heterogeneous
glioma dataset.

Future work may include extending the pipeline to multi-class tumor segmentation, exploring
additional preprocessing strategies such as histogram matching across patients, investigating
transformer-based segmentation models to further improve boundary accuracy, and comparing
patient-level against image-level cross-validation results across all seven configurations.


%% ── BIBLIOGRAPHY ─────────────────────────────────────────────────────────────
\clearpage
\addcontentsline{toc}{section}{References}
\printbibliography

\end{document}
